% bookedge-README.tex —— bookedge 模板使用说明（xelatex 编译）
\documentclass[zihao=5]{ctexart}
\usepackage[margin=2.6cm]{geometry}
\usepackage{xcolor}
\definecolor{inkred}{HTML}{A6433A}
\definecolor{inkgray}{HTML}{444444}
\usepackage{fancyhdr}
\pagestyle{fancy}\fancyhf{}\fancyfoot[C]{\small\color{inkgray}\thepage}
\renewcommand{\headrulewidth}{0pt}
\usepackage{listings}
\lstset{basicstyle=\ttfamily\small,breaklines=true,
  frame=single,rulecolor=\color{inkgray!40},backgroundcolor=\color{inkgray!4},
  commentstyle=\color{inkgray!60},language=TeX,showstringspaces=false}

\newcommand{\cmd}[1]{{\ttfamily\color{inkred}\detokenize{#1}}}

\title{\bfseries bookedge 书边索引书籍模板\\[2mm]\large 使用说明}
\author{}
\date{2026 年 11 月 · v1.0}

\begin{document}
\maketitle
\vspace{-6mm}

\section{这是什么}

\cmd{bookedge} 是一套 XeLaTeX 书籍模板，核心特征是\textbf{书口索引（thumb
index）}：每一章在页面外缘拥有一个专属颜色的色块，当前章的色块加宽加深，
其余章以浅色呈现。翻动书页时，拇指所在的书口位置正好压在这些色块上——
像词典一样，凭手感就能翻到想要的章。

配套的视觉系统：每章的专属色同时用于章号、节号、要点框、目录色签与页码
徽章，形成章内一致的色彩氛围；「部」页整页铺该部主色的极浅底并加一条横
色带，全书按部循环十色色板。

\section{文件}

\begin{tabular}{ll}
\cmd{bookedge.cls}      & 类文件（唯一必需）\\
\cmd{bookedge-demo.tex} & 示例书稿：一本 28 页的演示书，含全部功能 \\
\cmd{bookedge-README.pdf} & 本说明 \\
\end{tabular}

\section{与 \textsf{thumbs} 宏包的区别}

CTAN 上已有的 \textsf{thumbs} 宏包（v1.0q，2014）也能在页缘做拇指索引，
但两者的设计模型不同：\textsf{thumbs} 的页缘\textbf{每页只画当前对象的
一个色块}，块的垂直位置随序号逐级下移——用「位置」编码「你在哪」；
bookedge 则\textbf{同时画出全书每一章的色块}，排成一条固定竖梯，当前章
实色加宽、其余章浅色占位——用「整条阶梯」同时表达全书结构与当前位置。
合上书，拇指扫过书口即可感知各章的相对位置与厚度，这是词典拇指索引的本意。

其他差异：
\begin{itemize}
\item \textsf{thumbs} 需要在每个标记点手工写
  \cmd{\addthumb\{标题\}\{文字\}\{字色\}\{底色\}}，颜色、文字全部手工指定；
  bookedge 自动挂接 \cmd{\chapter}，章色按十色板循环，无需干预。
\item \textsf{thumbs} 的独有亮点是 \cmd{\thumbsoverview}——自动生成带
  超链接的「索引总览页」（bookedge 目前没有此功能，目录页的彩色章签
  承担类似职责）。
\item \textsf{thumbs} 是通用宏包，article/report 里任何对象都能加标记；
  bookedge 是书籍模板，索引条与章首导航尺、部页、封面是一套整体设计。
\item 兼容性实测：\textsf{thumbs} 与 ctexbook 同用需显式
  \cmd{\pagenumbering\{arabic\}}，否则其依赖的 \textsf{pageslts} 报错。
\end{itemize}

选型建议：给任意文档、任意位置加拇指标记或需要总览页，用 \textsf{thumbs}；
写中文书、要词典式书口导航并与整套版式咬合，用 bookedge。

\section{快速上手}

把 \cmd{bookedge.cls} 与你的正文放在同一目录，正文骨架如下：

\begin{lstlisting}
\documentclass{bookedge}          % 默认 openright（章起正面）；电子稿可用 openany
\bookedgetitle{书名}
\title{书名}
\epigraph{副书名}
\author{作者名\quad 编著}
\publishinfo{年份\quad 版次\quad 说明}
\begin{document}
\maketitle                        % 封面（右侧十色条自动绘制）
% 扉页、版权页：参看 bookedge-demo.tex 的写法
\chapter*{前言}
\tableofcontents
\bookedgetabson                   % 从目录页起启用书边索引
\mainmatter
\part{第一部分}
\chapter{第一章标题}
\begin{entry}{术语}{Foreign Term} % 词条框
术语解释正文……
\end{entry}
\begin{notebox}{要点}             % 要点框
要点内容……
\end{notebox}
正文……
\end{document}
\end{lstlisting}

用 \cmd{xelatex} 编译\textbf{两到三遍}至不再出现 Rerun 提示：书边索引的
章总数、\cmd{remember picture} 页面坐标都需要辅助文件回传，多编译一遍
即可收敛。

\section{主要接口}

\begin{tabular}{@{}p{5.2cm}p{9.2cm}@{}}
\cmd{\bookedgetitle\{书名\}}   & 设置书名（页眉反面、PDF 元信息）\\
\cmd{\epigraph\{副书名\}}      & 封面题词（书名下方）\\
\cmd{\publishinfo\{…\}}       & 版权页信息（版次等）\\
\cmd{\bookedgetabson}         & 启用书边索引（建议放在 \cmd{\tableofcontents} 前）\\
\cmd{\bookedgetotal\{n\}}     & 手工指定章总数（默认自动从辅助文件统计）\\
\cmd{\begin\{entry\}\{术语\}\{外文\}} & 词条框环境：章色标签 + 中外文对照，正文写环境体内 \\
\cmd{\begin\{notebox\}\{标题\}} & 要点框环境 \\
\cmd{\begin\{examplebox\}\{标题\}} & 示例框环境 \\
\end{tabular}

\section{定制}

\begin{itemize}
\item \textbf{换色板}：\cmd{bookedge.cls} 的「色板」一节有十条
  \cmd{\definecolor}（\cmd{beC1}$\sim$\cmd{beC10}，低饱和、可长读的莫兰迪
  系），改这里全书与索引条同步更新。
\item \textbf{开本}：默认 170\,mm$\times$240\,mm、版心 118\,mm，在导言区
  \cmd{\geometry} 处调整。
\item \textbf{索引条粗细}：\cmd{beWact}/\cmd{beWina}（当前章与其余章的
  色块宽度，单位 pt）。
\item \textbf{起页方式}：默认 \cmd{openright}——章与部一律起于正面
  （奇数页），必要时自动垫白，符合印刷书规范；纯屏幕阅读可加类选项
  \cmd{openany} 取消垫白。垫白页同样绘制索引条，保证书口阶梯连续。
\end{itemize}

\section{黑白打印：灰度版}

黑白打印时加类选项 \cmd{gray} 即可，正文无需任何改动：

\begin{lstlisting}
\documentclass[gray]{bookedge}    % 10 级灰阶 + 黑色定位标记
\end{lstlisting}

灰度版的取色逻辑与彩色版不同，专为单色印刷设计：

\begin{itemize}
\item 章节区分由\textbf{十级灰阶}承担（第 1 章最浅、第 10 章最深，
  相邻档约 8\%，黑白激光打印件上可分辨）；书口索引条、章首导航尺、
  目录色签、页码徽章、词条框标题条都按本章灰阶上色。
\item 当前章在索引条与导航尺上改用\textbf{纯黑加宽}标记——单色印刷
  下最强的“你在哪里”信号，不受灰阶深浅影响。
\item 深灰底自动配白字、浅灰底自动配深字，无需干预。
\item 正文小字（节号、项目符号等）统一用深灰而非本章灰阶，避免浅灰
  小字在纸上发虚。
\end{itemize}

\section{已知约定}

\begin{itemize}
\item 全书章数 $\le 10$ 时索引一色一章；超过 10 章色板自动循环，索引条
  会变薄，此时建议改以「部」分色（部色同样循环）。
\item 封面、扉页、版权页不显示索引；从目录页起，索引以全浅色状态出现，
  相当于全书的图例。
\item 附录等 \cmd{\chapter*} 页沿用前一章的颜色与索引位置。
\item 索引条与页眉的左右取向按\textbf{物理页序}判断，即使正文重置了页
  码计数（如 \cmd{\mainmatter}），色块也始终落在正确的书口一侧。
\end{itemize}

\end{document}
